\chapter*{Preface}
\markboth{Preface}{Preface}
\addcontentsline{toc}{chapter}{Preface}
These notes are for students who know the main ideas of single-variable
calculus and want to understand why its central statements are true.  They
develop rigorous real analysis before rebuilding calculus from limits,
completeness, continuity, differentiation, and Riemann integration.  The
intended audience is an advanced undergraduate beginning proof-based
mathematics. Familiar algebra and ordinary calculus are assumed, but the
set-theoretic, logical, topological, and finite-dimensional linear-algebraic
tools needed later are developed in the text.

Part~I (Chapters~1--9) constructs the real number system and develops
one-variable analysis through Riemann integration and the Fundamental
Theorem of Calculus. Part~II (Chapters~10--13) develops the finite-dimensional
topology and linear algebra needed for multivariable differentiation and
Riemann integration, and ends with differential forms and the general Stokes
theorem. Measure theory, Lebesgue integration, complex and functional
analysis, partial differential equations, and deeper differential geometry
are outside the developed theory; Chapter~9 previews the Lebesgue FTC,
and the final chapter describes further directions.

\paragraph{AI-assisted preparation.}
OpenAI Codex was used extensively in drafting, revising, proofreading, and
preparing the \LaTeX{} source for these notes. The author determined the
scope and mathematical organization, made the final decisions about retained
content, and is responsible for approving and verifying the correctness of
the published version.

\chapter*{How to Read These Notes}
\markboth{How to Read These Notes}{How to Read These Notes}
\addcontentsline{toc}{chapter}{How to Read These Notes}
These notes start from familiar calculus questions but use a different
standard of explanation: every formal assertion has hypotheses, and every
main claim is justified from material already established or from an
explicitly identified external theorem. A definition fixes meaning; a theorem
requires proof unless it is marked as external. Examples provide intuition,
but not proof of a general assertion. Occasionally, an explicitly identified
informal preview uses familiar notation to motivate a later formal
development; such previews are not used as formal justification.

\paragraph{Practice and hints.}
Part~I gives explicit proof strategies, ordinary calculations, and a few
applications that require translating words into mathematical assumptions.
The scaffolding becomes lighter in Part~II as the reader gains proof
maturity; its structural development expects greater independence.
An asterisk \textsuperscript{*} after an exercise number indicates that
a hint is available in Appendix~\ref{app:hints}; the symbol links directly
to that hint, and the hint's exercise number links back. The asterisk means
\emph{hint available}, not ``difficult exercise.'' Some theoretical exercises
are deliberately left without hints because finding the idea is part of
their purpose.

\paragraph{Notation and reminders.}
General functions and nonlinear maps use lowercase Latin letters; charts
use lowercase Greek letters, curves use $\gamma$, and surface
parametrizations usually use $\sigma$. Linear maps use capitals such as
$T,S,L$. The Euclidean derivative is $Df(a)$; the intrinsic differential
and exterior derivative are $df_p$ and $d\omega$. For a smooth Euclidean
scalar function, $df$ denotes the exact $1$-form and $(df)_p=Df(p)$
its value at $p$. Conventional symbols
with a distinct role, such as the force or flux vector field $F$ and an
antiderivative $F$ paired with its integrand $f$, are retained where they
clarify that role. Less familiar structural notions reused after a gap
receive a cross-reference; when central to the argument, they also receive
a local reminder. Such reminders are selective rather than attached to
every use of ordinary continuity or differentiation.

The chapters are designed to be read in order. Exercises are part of the
course: they supply routine practice, short proofs, counterexamples, and
connections among the main results. A reader should attempt them after the
relevant section, consulting the stated hypotheses and earlier results rather
than relying on calculus folklore. Results stated without proof are identified where they occur. The
bibliography also distinguishes the imported topological result from
further-reading references. In particular, Chapter~9 states without proof
the Lipschitz/a.e. and Lebesgue versions of the Fundamental Theorem of
Calculus and the existence of the Cantor function. These previews clarify
the limits of the Riemann theory but are not prerequisites for later chapters.
The optional analytic Lagrange inversion result in Chapter~11 and the
Lebesgue change-of-variables preview in Chapter~12 are likewise stated
without proof. Section~13.17 contains additional explicitly unproved preview theorems,
none used as a prerequisite. The main route through these chapters is local differential
structure, primitive substitution with finite Euclidean localization, and
then locally finite manifold localization in Chapter~13.


Read a substantial proof twice with different aims. On the first pass,
identify the problem, the main construction, and the purpose of each estimate.
On the second, reconstruct the quantifiers and check where every hypothesis
enters. Understanding a proof's architecture and reproducing all its details
are related but distinct skills.

In Chapter~\ref{ch:reals}, the construction of $\R$ proves that the
system we use exists. Difficulty reproducing
every equivalence-class verification should not block further study.
Retain its working consequences: ordered-field structure, the Archimedean
property, rational density, completeness, Cauchy convergence, and the
least-upper-bound property. Return to the implementation as these ideas
become familiar. In Part~II, watch these same tools reappear as norm estimates,
compactness arguments, linear approximation, and local-to-global integration.

Chapter~10 assumes no previous topology or linear algebra course. Its
distance examples lead to open sets and metric continuity before abstract
topology; its coordinate calculations lead to bases, direct sums, linear
maps, and matrices. Appendices~\ref{app:foundations} and~\ref{app:categories}
collect foundational and structural language occasionally useful throughout
the notes. Appendices~\ref{app:determinants}--\ref{app:cohomology} develop
optional algebraic, topological, differential-geometric, and cohomological
extensions. None is intended to be read before Chapter~1, and none is a
prerequisite for the main analysis-to-Stokes route unless explicitly stated.
The hierarchy is foundational language, then categorical language, then
optional mathematical extensions; it is not a compulsory reading order.

Appendix~\ref{app:foundations} distinguishes induction, recursion, and choice.
Appendix~\ref{app:categories} organizes constructions through maps and their
compatibility. Appendix~\ref{app:determinants} develops determinants over a
field; Appendix~\ref{app:tensors} gives a structural alternative to
Chapter~13's concrete exterior algebra. Appendix~\ref{app:quotients}
constructs manifolds by gluing. Appendix~\ref{app:submanifolds} connects
the rank theorems to manifold geometry, and Appendix~\ref{app:lie} uses
that framework for Lie groups. Appendix~\ref{app:cohomology} starts from
forms and Stokes, preserves concrete computations, and then introduces
the exact-sequence language connecting smooth analysis with topology.
These are branching continuations, not one total dependency chain.
Section~13.17 provides navigation, including differential topology,
geometric topology, and geometric analysis; the appendices develop selected
parts of that outlook.

At the first serious use of a proof method, look for the actual map,
sequence, or estimate being constructed. Later uses refer back to that
mechanism. In particular, existence of a map is a different task from
proving equality of two maps, and choosing a finite index is a different
task from subsequently taking a limit.
